\subsection{Incremental principal component analysis}
Conventional Principal component analysis requires storing all observations made over time and calculating full matrix eignedecomposition, introducing a high computational complexity which is unsuitable for real time anomaly detection. A smart solution is to shift the computational load to functions which are able to reuse already calculated values for its computation lowering the computational cost massively.Following the incremental principal component analysis framework introduced by \textcite{Incremental_PCA}, the proposed pipeline maintains running data statistics and updates the dominant eigenspace iteratively upon receiving each new observation.
Even with this optimization steps a central challenge remains: extracting the eigenvectors of the covariance matrix at every iteration step.  
According to \textcite[p.~2931]{1468483}, note that traditional signal subspace estimation relies on eigenvalue or singular value decomposition, whose inherent computational complexity creates a genuine need for fast subspace tracking techniques in adaptive signal processing. This motivates the subspace tracking approach, which instead avoids full eignedecomposition by iteratively refining an estimate of the dominant eigenvectors.